\section{Uniform lower approximation and closed thresholds}\label{sec:hierarchy}
For a channel $\Phi:A\to B$, choose a complementary channel $\Phi^c$ and put
\[
 b_n(\Phi)=\max_{\rho} \{S(\Phi^{\otimes n}(\rho))
                  -S((\Phi^c)^{\otimes n}(\rho))\},\qquad
 a_n(\Phi)=b_n(\Phi)/n.
\]
All logarithms are base two unless $\ln$ is written. Pure inputs give
$b_n\geq0$. Product inputs give $b_{n+m}\geq b_n+b_m$, so
$b_n$ is nondecreasing, although $a_n$ need not be. The capacity formula
and Fekete's lemma give
\begin{equation}\label{eq:capacity}
 Q(\Phi)=\lim_{n\to\infty}a_n(\Phi)=\sup_{n\geq1}a_n(\Phi).
\end{equation}

\begin{theorem}[Uniform finite-block and capacity approximation]\label{thm:uniform}
There is a uniform algorithm returning rationals $c(\Phi,n,k)$ such that
\begin{equation}\label{eq:c}
 |c(\Phi,n,k)-a_n(\Phi)|\leq2^{-k}.
\end{equation}
Consequently $Q$ is uniformly left-c.e.: a computable increasing rational
sequence $r_s(\Phi)$ converges to $Q(\Phi)$.
\end{theorem}
\begin{proof}
Here is the finite-block algorithm; Appendix~\ref{app:effective} supplies
two explicit effective-net derivations. A reference system of dimension
$D=d_A^n$ purifies every input. Thus $b_n$ is the maximum of
$-S(R|B)$ over a unit sphere in $\C^D\otimes\C^D$.
For unit vectors at Euclidean distance at most $\eta\leq1$, their output
states have trace distance at most $\eta$. Winter's bound gives the
computable modulus
\begin{equation}\label{eq:modulus}
 |I_c(\psi)-I_c(\varphi)|\leq
 2\eta n\log d_A+g(\eta),\quad
 g(\eta)=(1+\eta)h_2\!\left(\frac{\eta}{1+\eta}\right),
\end{equation}
where $g(0)=0$. Optimize certified entropy evaluations over a sufficiently
fine finite net. Both the mesh error and evaluation error are controlled
uniformly. Divide by $n$ to obtain~\eqref{eq:c}.

The following lower-approximation step is the effective Fekete principle
\cite[Lemmas~17--19 and Theorem~21]{BBD}, applied uniformly to the
finite-block channel quantities. Explicitly, define
\begin{equation}\label{eq:rs}
 r_s=\max\bigl(\{0\}\cup
 \{c(\Phi,n,k)-2^{-k}:1\leq n,k\leq s\}\bigr).
\end{equation}
These are increasing rational lower bounds. Given any $a_n$, taking
$k$ arbitrarily large shows $\sup_s r_s\geq a_n$.
Equation~\eqref{eq:capacity} proves convergence to $Q$.
\end{proof}

\begin{corollary}[Threshold bookkeeping]\label{cor:threshold}
For every fixed computable threshold $t$, the predicate $Q(\Phi)>t$ is
c.e., $Q(\Phi)\geq t$ is in $\Pi^0_2$, and $Q(\Phi)<t$ is in
$\Sigma^0_2$. In particular, (QCap) is in $\Pi^0_2$,
(GapQcap) is in RE, and $Q=0$ is co-RE.
\end{corollary}
\begin{proof}
For rational $t$ the decisive formulas have decidable matrices:
\begin{align}
 Q>t&\ \Longleftrightarrow\ \exists s\ [r_s>t],\label{eq:strict}\\
 Q\geq t&\ \Longleftrightarrow\
 \forall k\geq1\ \exists s\ [r_s>t-2^{-k}],\label{eq:closed}\\
 Q<t&\ \Longleftrightarrow\
 \exists k\geq1\ \forall s\ [r_s\leq t-2^{-k}].\label{eq:below}
\end{align}
For computable $t$, strict comparison with it is c.e.; its precision
witness merges with the existential stage in~\eqref{eq:closed}.
Complementation gives the last threshold class. Search for $r_s>1/2$
accepts every promised $Q\geq3/4$ instance and no promised $Q\leq1/4$
instance. Finally $Q\geq0$, so~\eqref{eq:strict} at zero proves the
positivity and zero-capacity claims.
\end{proof}
The RE upper bound for (GapQcap) is already stated by
BMZ~\cite[Section~1.3]{BMZ}; we include it to compare the gapped and
closed thresholds.

Equality at $Q=t$ is included in~\eqref{eq:closed}: every threshold
$t-2^{-k}$ is eventually crossed. This is why numerical approximation
does not turn the closed threshold into an existential test. For general
program names one can dovetail completed lower evaluations, including
partial computations off the valid-input promise. This gives a
$\Pi^0_2$ extension agreeing on valid names. No validity test for such
names is required.

\paragraph{BMZ's conjecture.}
Computable preimages preserve $\Pi^0_2$. Thus $\Pi^0_3$ many-one hardness
of (QCap) would collapse the strict arithmetical hierarchy. Under Turing
reductions, a $\Pi^0_2$ language is computable in $\mathbf0''$, whereas
a $\Pi^0_3$-complete language has degree $\mathbf0'''$. The same
contradiction holds for valid-query promise reductions, using the
extension just described. BMZ's stated $\Pi^0_4$ membership remains a
valid loose bound. Their separate gap-hardness conjecture is unaffected.

\begin{conjecture}[A corrected form]\label{conj:qcap}
On the effective algebraic circuit domain, (QCap) is $\Pi^0_2$-complete.
\end{conjecture}
This is a proposed replacement, not a lower bound proved here.

\paragraph{Question 2 and finite attainment.}
BMZ ask whether a channel exists with $a_n(\Phi)<Q(\Phi)$ for every
integer $n\geq1$. The following implication uses those integers, rather
than only powers of two.
\begin{proposition}\label{prop:attainment}
If every channel in the encoded family attains $Q$ at some finite block
length, then (QCap) is in $\Delta^0_2$. Hence $\Pi^0_2$ hardness, under
many-one or Turing reductions, would force a channel in that family
with $a_n<Q$ for all $n$.
\end{proposition}
\begin{proof}
Under finite attainment, $Q\geq t$ is equivalent to
$\exists n\ [a_n\geq t]$. For rational $t$,
\[
 a_n\geq t\ \Longleftrightarrow\
 \forall k\ [c(\Phi,n,k)+2^{-k}\geq t].
\]
Thus the threshold is $\Sigma^0_2$ and, independently by
Corollary~\ref{cor:threshold}, $\Pi^0_2$. Their intersection is
$\Delta^0_2$, computable in $\mathbf0'$. This excludes $\Pi^0_2$
Turing hardness as well as many-one hardness.
\end{proof}
No all-length nonattaining channel follows without the hardness
hypothesis. Nor does undecidability of an exact threshold alone imply
noncomputability of the capacity function: equality tests for computable
reals can already be undecidable. By contrast, many-one RE-hardness of
this closed threshold would exclude a uniformly computable $Q$, since
its closed threshold would then be co-c.e.; RE-hardness of the gapped
promise would also exclude it, since a computable $Q$ decides that promise.
